\subsection{Graph of a Linear Function} 
\begin{theorem}
Both of the following are true.
\begin{itemize}
\item If a function is linear, then its graph is a straight line.
\item If a graph is a straight line (not horizontal or vertical), it is the graph of a linear function.
\end{itemize}
\end{theorem}
\begin{example}
Fill in the table below for $f(x)=\frac{2}{3}x-4$. Plot each pair in the table, then sketch the graph of $f$.
\begin{lpic}{}{10}
\regtext{1}{-2}{$x$};
\regtext{4}{-2}{$y=f(x)$};
\regtext{1}{-3}{$-3$};
\foreach \y in {0,3,6}
	\regtext{1}{-4-\y/3}{\figdash$\y$};
\draw (0,-2) -- (6,-2);
\draw (2,-1) -- (2,-6);
\begin{scope}[xshift=\value{cols}*2\linespace-8\linespace,yshift=-2\linespace]
\setgraph{4}{7}{7}{1}
\setspace{1}
\AxesSmall
\end{scope}
\end{lpic}
\end{example}
\begin{note}
There are two types of lines that are not graphs of linear functions.
\begin{itemize}
\item \term{Horizontal} lines have equations in the form $\makeblanksmall=\mbox{a number}$. We say this is a \makeblank[1.5in] function.
\item \term{Vertical} lines have equations in the form $\makeblanksmall=\mbox{a number}$. This graph fails the \makeblank, it is \makeblank[1in] the graph of a function.
\end{itemize}
The lines that graph linear functions are \term{oblique} lines.
\end{note}
